% =============================================================================
% APENDICE: PARAMETROS OPTIMOS DE UMBRAL
% =============================================================================

\section{Par\'ametros \'Optimos de Umbral}
\label{app:optimal_params}

Los umbrales de posicionamiento $(q_{\text{extreme}}, q_{\text{moderate}})$ se determinan din\'amicamente para cada modelo y cada d\'ia del per\'iodo de prueba mediante un metamodelo \textit{K-Nearest Neighbors} (KNN) entrenado exclusivamente con datos del per\'iodo de entrenamiento. A continuaci\'on se detallan los par\'ametros del procedimiento y las configuraciones resultantes.

\subsection{Metodolog\'ia de Optimizaci\'on mediante Meta-KNN}
\label{app:opt_methodology}

La descripci\'on narrativa completa del procedimiento, incluyendo la construcci\'on del meta-dataset y las siete \textit{meta-features} que describen el r\'egimen predictivo de cada ventana, se presenta en la Secci\'on~\ref{subsec:meta_knn}; este ap\'endice documenta la especificaci\'on exacta empleada. El espacio de b\'usqueda comprende las 26 combinaciones v\'alidas con $q_{\text{extreme}} \in \{5, 10, 15, 20, 25, 30\}$, $q_{\text{moderate}} \in \{20, 30, 40, 50, 60\}$ y la restricci\'on $q_{\text{moderate}} > q_{\text{extreme}}$, dado que la posici\'on apalancada exige mayor convicci\'on que la moderada. El metamodelo emplea $K = 5$ vecinos ponderados por distancia inversa sobre \textit{meta-features} estandarizadas mediante \textit{StandardScaler}, con ventanas deslizantes de 63 d\'ias que avanzan en pasos de 21 y un horizonte de evaluaci\'on de otros 63 d\'ias, lo que produce alrededor de 242 muestras de entrenamiento por modelo.

Para cada ventana de evaluaci\'on del meta-dataset se selecciona la combinaci\'on que maximiza el \textit{Sharpe ratio} penalizado.

\begin{equation}
\text{score}(q_{\text{ext}}, q_{\text{mod}}) =
\begin{cases}
\text{Sharpe}(R_{\text{net}}) - 0{,}5 & \text{si } \text{MaxDD} > 20\% \\[4pt]
\text{Sharpe}(R_{\text{net}}) & \text{en caso contrario}
\end{cases}
\end{equation}

\noindent La penalizaci\'on de $0{,}5$ desincentiva combinaciones que generan \textit{drawdowns} excesivos. Se aplica adem\'as un filtro basado en la desviaci\'on est\'andar \textit{rolling} de las predicciones; si $\sigma_{\text{rolling},63} < \tau = 0{,}001$, la estrategia se posiciona en efectivo con independencia de los umbrales predichos por el KNN, lo que protege contra predicciones cuasiconstantes que generar\'ian percentiles aleatorios.

\subsection{Tabla de Par\'ametros Din\'amicos}
\label{app:optimal_table}

La Tabla~\ref{tab:optimal_params} presenta los par\'ametros din\'amicos m\'as frecuentes para cada modelo.

\begin{table}[htbp]
\centering
\caption{Par\'ametros \'Optimos de Umbral por Modelo}
\label{tab:optimal_params}
\small
\begin{tabular}{@{}llrrr@{}}
\toprule
\textbf{Modelo} & \textbf{Categor\'ia} & \textbf{$q_{3x}$ (\%)} & \textbf{$q_{1x}$ (\%)} & \textbf{Interpretaci\'on} \\
\midrule
LSTM\_Attention & Deep learning & 5 & 50 & Top 5\% $\rightarrow$ 3x, Top 50\% $\rightarrow$ 1x \\
Ridge & ML cl\'asico & 20 & 30 & Top 20\% $\rightarrow$ 3x, Top 30\% $\rightarrow$ 1x \\
RandomForest & ML cl\'asico & 5 & 40 & Top 5\% $\rightarrow$ 3x, Top 40\% $\rightarrow$ 1x \\
CNN\_LSTM & Deep learning & 20 & 40 & Top 20\% $\rightarrow$ 3x, Top 40\% $\rightarrow$ 1x \\
DLinear & Deep learning & 5 & 60 & Top 5\% $\rightarrow$ 3x, Top 60\% $\rightarrow$ 1x \\
GradientBoosting & Boosting & 5 & 20 & Top 5\% $\rightarrow$ 3x, Top 20\% $\rightarrow$ 1x \\
CatBoost & Boosting & 5 & 20 & Top 5\% $\rightarrow$ 3x, Top 20\% $\rightarrow$ 1x \\
Theta & Series temporales & 15 & 60 & Top 15\% $\rightarrow$ 3x, Top 60\% $\rightarrow$ 1x \\
NBEATS & Deep learning & 10 & 20 & Top 10\% $\rightarrow$ 3x, Top 20\% $\rightarrow$ 1x \\
\rowcolor{naivebg} SeasonalNaive$^\dagger$ & Series temporales & 5 & 60 & Top 5\% $\rightarrow$ 3x, Top 60\% $\rightarrow$ 1x \\
Lasso & ML cl\'asico & 5 & 40 & Top 5\% $\rightarrow$ 3x, Top 40\% $\rightarrow$ 1x \\
GARCH & Series temporales & 5 & 30 & Top 5\% $\rightarrow$ 3x, Top 30\% $\rightarrow$ 1x \\
LightGBM & Boosting & 5 & 20 & Top 5\% $\rightarrow$ 3x, Top 20\% $\rightarrow$ 1x \\
NHiTS & Deep learning & 5 & 60 & Top 5\% $\rightarrow$ 3x, Top 60\% $\rightarrow$ 1x \\
AutoARIMA & Series temporales & 10 & 30 & Top 10\% $\rightarrow$ 3x, Top 30\% $\rightarrow$ 1x \\
TCN & Deep learning & 5 & 30 & Top 5\% $\rightarrow$ 3x, Top 30\% $\rightarrow$ 1x \\
TFT & Deep learning & 20 & 60 & Top 20\% $\rightarrow$ 3x, Top 60\% $\rightarrow$ 1x \\
Prophet & Series temporales & 15 & 60 & Top 15\% $\rightarrow$ 3x, Top 60\% $\rightarrow$ 1x \\
ElasticNet & ML cl\'asico & 10 & 60 & Top 10\% $\rightarrow$ 3x, Top 60\% $\rightarrow$ 1x \\
XGBoost & Boosting & 5 & 20 & Top 5\% $\rightarrow$ 3x, Top 20\% $\rightarrow$ 1x \\
BiLSTM & Deep learning & 5 & 20 & Top 5\% $\rightarrow$ 3x, Top 20\% $\rightarrow$ 1x \\
ExponentialSmoothing & Series temporales & 30 & 60 & Top 30\% $\rightarrow$ 3x, Top 60\% $\rightarrow$ 1x \\
BiGRU & Deep learning & 30 & 40 & Top 30\% $\rightarrow$ 3x, Top 40\% $\rightarrow$ 1x \\
\bottomrule
\end{tabular}
\begin{tablenotes}
\small
\item Nota: $q_{3x}$ = percentil para posici\'on UPRO (3x). $q_{1x}$ = percentil para posici\'on SPY (1x). Predicciones por debajo del percentil $q_{1x}$ resultan en posici\'on Cash (0). \colorbox{naivebg}{$^\dagger$SeasonalNaive} = \textit{benchmark} de complejidad m\'inima. Par\'ametros mostrados corresponden a la combinaci\'on m\'as frecuente seleccionada por el Meta-KNN din\'amico; los umbrales reales var\'ian por d\'ia.
\end{tablenotes}
\end{table}


\subsection{Ejemplos Ilustrativos de la Estrategia}
\label{sec:appendix_strategy_examples}

\paragraph{Ejemplo de umbrales din\'amicos con Ridge.} Para ilustrar c\'omo los umbrales se adaptan al contexto, consid\'erese el modelo Ridge durante cinco d\'ias consecutivos. En $t$, $t{+}1$ y $t{+}2$, el mercado presenta baja volatilidad; el Meta-KNN identifica este patr\'on como similar a ventanas hist\'oricas donde los umbrales $(q_{\text{extreme}}{=}20, q_{\text{moderate}}{=}30)$ resultaron \'optimos, es decir, UPRO si el percentil $\geq 80$ y SPY si $\geq 70$. En $t{+}3$, un incremento en la volatilidad del mercado altera las \textit{meta-features} y el KNN selecciona $(10, 40)$, lo que implica UPRO si $\geq 90$ y SPY si $\geq 60$. La Tabla~\ref{tab:ejemplo_posiciones} ilustra c\'omo un mismo percentil puede generar posiciones distintas seg\'un los umbrales vigentes.

\begin{table}[H]
\centering
\caption{Ejemplo de Conversi\'on de Predicciones a Posiciones con Umbrales Din\'amicos (Ridge)}
\label{tab:ejemplo_posiciones}
\small
\begin{tabular}{@{}crrccl@{}}
\toprule
\textbf{D\'ia} & $\hat{y}_t$ & \textbf{Pctil} & $(q_{\text{ext}}, q_{\text{mod}})$ & \textbf{Umbrales} & \textbf{Posici\'on} \\
\midrule
$t$ & $-0.005$ & 45 & (20, 30) & $\geq$80 / $\geq$70 & Cash \\
$t{+}1$ & $+0.002$ & 74 & (20, 30) & $\geq$80 / $\geq$70 & SPY (+1) \\
$t{+}2$ & $+0.004$ & 91 & (20, 30) & $\geq$80 / $\geq$70 & UPRO (+3) \\
$t{+}3$ & $+0.003$ & 84 & (10, 40) & $\geq$90 / $\geq$60 & SPY (+1) \\
$t{+}4$ & $-0.000$ & 53 & (10, 40) & $\geq$90 / $\geq$60 & Cash \\
\bottomrule
\end{tabular}
\end{table}

N\'otese que la predicci\'on de $t{+}2$ (percentil~91) activa la posici\'on apalancada bajo los umbrales $(20, 30)$, pero en $t{+}3$ un percentil de~84 (que habr\'ia activado UPRO bajo esos mismos umbrales) genera solo posici\'on SPY porque el Meta-KNN ha elevado el requisito a percentil~$\geq 90$ al detectar un cambio en las condiciones de mercado. Este mecanismo adapta autom\'aticamente no solo la escala de las predicciones (funci\'on del percentil \textit{rolling}) sino tambi\'en la agresividad del posicionamiento (funci\'on de los umbrales din\'amicos).

La Tabla~\ref{tab:optimal_params_winners} resume las caracter\'isticas de los cuatro modelos que superan al \textit{benchmark} SPY Buy \& Hold (+101\%, Sharpe 0{,}66) con la metodolog\'ia Meta-KNN. La diversidad en el n\'umero de combinaciones empleadas refleja la naturaleza din\'amica del sistema, dado que LSTM+Attention utiliza 22 configuraciones distintas de umbrales a lo largo de los 1{,}302 d\'ias de prueba, con su combinaci\'on m\'as frecuente representando apenas el 26{,}5\% de los d\'ias, mientras que RandomForest mantiene una configuraci\'on dominante el 83{,}9\% del tiempo. Ninguno de estos cuatro modelos se ve afectado significativamente por el filtro de calidad, puesto que todos superan $\tau$ al menos el 53\% de los d\'ias, lo que sugiere que sus predicciones contienen variaci\'on susceptible de ser explotada por el sistema de percentiles.

\paragraph{Ejemplo de \textit{backtesting} con costos de transacci\'on.} Para ilustrar la interacci\'on entre posiciones, retornos y costos, la Tabla~\ref{tab:backtest_example} presenta cinco d\'ias consecutivos hipot\'eticos del modelo LSTM+Attention. Se supone que el Meta-KNN selecciona la configuraci\'on $(q_{\text{extreme}} = 5, q_{\text{moderate}} = 20)$ durante estos d\'ias, lo que implica que la posici\'on UPRO se activa cuando el percentil alcanza o supera~95 y la posici\'on SPY cuando alcanza o supera~80. Se parte de un capital inicial de \$10{,}000 y una tasa libre de riesgo diaria de $r^f = 0{,}02\%$ ($\approx$5\% anual).

\begin{table}[H]
\centering
\caption{Ejemplo de \textit{Backtesting}, Cinco D\'ias con LSTM+Attention}
\label{tab:backtest_example}
\small
\begin{tabular}{@{}crcrlrrr@{}}
\toprule
\textbf{D\'ia} & \textbf{Pctil} & \textbf{Pos.} & $r^{SPY}_t$ & \textbf{Cambio} & \textbf{Costos} & $r^{estrat}_t$ & \textbf{Capital} \\
\midrule
$t$ & 62 & Cash & +0.80\% & --- & 0 bps & +0.02\% & \$10{,}002 \\
$t{+}1$ & 83 & SPY & +1.20\% & Cash$\to$SPY & 2 bps & +1.18\% & \$10{,}120 \\
$t{+}2$ & 96 & UPRO & +0.50\% & SPY$\to$UPRO & 10 bps & +1.36\% & \$10{,}257 \\
$t{+}3$ & 97 & UPRO & $-$0.30\% & --- & 3 bps & $-$0.97\% & \$10{,}157 \\
$t{+}4$ & 71 & Cash & +0.40\% & UPRO$\to$Cash & 5 bps & $-$0.03\% & \$10{,}154 \\
\bottomrule
\end{tabular}
\vspace{0.2cm}
\parbox{0.95\textwidth}{\footnotesize\textit{Nota: La columna Costos muestra los costos totales diarios redondeados a puntos b\'asicos. En $t$, la estrategia est\'a en Cash y gana solo $r^f$, evitando el +0.80\% del mercado. En $t{+}2$, la posici\'on 3x amplifica el exceso de retorno: $3 \times (0.50\% - 0.02\%) = +1.44\%$, m\'as $r^f$ da +1.46\% bruto; se restan 10~bps de costos: 7~bps de \textit{bid-ask} (salida de SPY + entrada a UPRO) y 3~bps de \textit{volatility drag}. En $t{+}3$, sin cambio de posici\'on, los costos se reducen a 3~bps diarios de \textit{volatility drag}; la amplificaci\'on 3x opera en contra: $3 \times (-0.32\%) = -0.96\%$ de exceso, produciendo $-$0.94\% bruto. En $t{+}4$, la salida de UPRO cuesta 5~bps de \textit{spread}, que exceden la tasa libre de riesgo (+0.02\%), resultando en un retorno neto ligeramente negativo ($-$0.03\%).}}
\end{table}

Este ejemplo ilustra un principio central de la estrategia, a saber, la selectividad temporal. El modelo LSTM+Attention permanece en efectivo el 55{,}1\% del tiempo, buscando capturar los momentos donde las \textit{meta-features} y las condiciones de mercado se asemejan a per\'iodos hist\'oricos de alto rendimiento. El costo de esta selectividad (perder d\'ias alcistas como $t$ y $t{+}4$) se compensa con la evitaci\'on de d\'ias bajistas y la amplificaci\'on 3x en los d\'ias de alta convicci\'on correcta. El an\'alisis completo del impacto de costos sobre los 23 modelos se presenta en la Secci\'on~\ref{sec:costs}.

\subsection{An\'alisis de Par\'ametros}
\label{app:param_analysis}

\subsubsection{Diversidad de Umbrales por Modelo}

La naturaleza din\'amica del Meta-KNN implica que cada modelo utiliza m\'ultiples combinaciones de umbrales a lo largo del per\'iodo de prueba. Se distinguen tres patrones.

\begin{enumerate}
    \item \textbf{Modelos altamente din\'amicos} (>15 combinaciones), entre los que destacan LSTM+Attention (22), Ridge (21) y CNN-LSTM (18). El KNN adapta activamente los umbrales a las condiciones cambiantes. Predominan las arquitecturas de \textit{deep learning} y modelos lineales con predicciones ricas.

    \item \textbf{Modelos moderadamente din\'amicos} (5--15 combinaciones), entre los que destacan CatBoost (10), Theta (9), SeasonalNaive (7), LightGBM (7) y GARCH (7). El KNN encuentra un rango limitado de configuraciones \'utiles.

    \item \textbf{Modelos est\'aticos o filtrados} (<5 combinaciones o mayoritariamente en Cash), entre los que destacan RandomForest (4), XGBoost (4), NBEATS (2), NHiTS (1), GradientBoosting (1) y AutoARIMA (0, 100\% filtrado). Predicciones insuficientemente variables o patr\'on dominante claro.
\end{enumerate}

\subsubsection{Implicaciones}

\begin{itemize}
    \item \textbf{Alta diversidad + buen rendimiento.} El modelo captura distintos reg\'imenes de mercado y el KNN adapta los umbrales adecuadamente (LSTM+Attention con 22 combinaciones y +395\%).
    \item \textbf{Baja diversidad + buen rendimiento.} Una configuraci\'on robusta funciona consistentemente (RandomForest con 4 combinaciones y +128\%).
    \item \textbf{Filtrado mayoritario.} Predicciones cuasiconstantes impiden la generaci\'on de se\~nales, y el filtro de calidad protege correctamente (AutoARIMA permanece 100\% en Cash).
\end{itemize}

\subsection{Interpretaci\'on de los Umbrales Din\'amicos}
\label{app:threshold_interpretation}

\subsubsection{Significado de los Percentiles}

Para LSTM+Attention, cuya combinaci\'on m\'as frecuente es $(q_{\text{extreme}}{=}5, q_{\text{moderate}}{=}50)$.

\begin{itemize}
    \item \textbf{Predicci\'on en top 5\%.} Posici\'on UPRO (3x \textit{leverage})
    \item \textbf{Predicci\'on en top 5--50\%.} Posici\'on SPY (1x)
    \item \textbf{Predicci\'on debajo del top 50\%.} Cash (tasa libre de riesgo)
\end{itemize}

Sin embargo, estos umbrales var\'ian din\'amicamente. En per\'iodos de alta volatilidad, el KNN puede seleccionar $(5, 20)$, reservando SPY solo para el top 20\%, mientras que en per\'iodos estables puede ampliar a $(20, 60)$, adoptando mayor exposici\'on. El modelo utiliz\'o 22 combinaciones distintas a lo largo del per\'iodo de prueba, lo que indica una adaptaci\'on activa a las condiciones de mercado.

\subsubsection{Percentiles vs Valores Absolutos}

El uso de percentiles \textit{rolling} en lugar de umbrales fijos tiene ventajas.

\begin{enumerate}
    \item \textbf{Adaptabilidad.} Se ajusta autom\'aticamente a cambios en la distribuci\'on de predicciones.
    \item \textbf{Robustez.} Menos sensible a \textit{drift} en el modelo.
    \item \textbf{Normalizaci\'on.} Permite comparar la convicci\'on de modelos con escalas muy distintas.
\end{enumerate}

Combinados con los umbrales din\'amicos del Meta-KNN, los percentiles \textit{rolling} proporcionan un sistema de posicionamiento doblemente adaptativo, tanto a la escala del modelo (v\'ia percentiles) como a las condiciones de mercado (v\'ia KNN).
